This index collects the notation used in Chapter 0, grouped by mathematical context. The links in the final column lead to the relevant definition or convention. Local meanings are recorded separately when the same letter serves different purposes.

\subsection*{Matrices and traces}\label{v1a-notation-matrices}

\begin{longtable}{p{0.22\linewidth}p{0.58\linewidth}p{0.08\linewidth}}
\textbf{Symbol} & \textbf{Meaning and convention} & \textbf{See} \\ \hline
\(A=(a_{ij})\) & Complex matrix, with entry \(a_{ij}\) in row \(i\), column \(j\). & \hyperref[definition-the-trace]{0.1} \\
$A^{\mathsf{T}}$ & Transpose: $(A^{\mathsf{T}})_{ij}=a_{ji}$; entries are not conjugated. A transposed row vector is a column vector. & \hyperref[linear-algebra.transpose]{0.1} \\
\(I_n\) & Identity matrix of size \(n\); its trace is \(n\). & \hyperref[definition-the-trace]{0.1} \\
\(M_n(\mathbb C)\) & Complex \(n\times n\) matrices. & \hyperref[linear-algebra.trace-similarity]{0.1} \\
\(\det(tI_n-A)\) & Characteristic polynomial of \(A\). & \hyperref[linear-algebra.trace-eigenvalue-sum]{0.1} \\
\(\lambda_j\) & Eigenvalues, repeated according to algebraic multiplicity. & \hyperref[linear-algebra.trace-eigenvalue-sum]{0.1} \\
\(\operatorname{tr}(A)\) & Sum of the diagonal entries; the empty sum is zero. & \hyperref[definition-the-trace]{0.1} \\
\(\operatorname{tr}(T)\) & Trace of a finite-dimensional endomorphism, computed in any basis. & \hyperref[linear-algebra.endomorphism-trace]{0.1} \\
\end{longtable}

\subsection*{Function spaces and finite kernels}\label{v1a-notation-kernels}

\begin{longtable}{p{0.22\linewidth}p{0.58\linewidth}p{0.08\linewidth}}
\textbf{Symbol} & \textbf{Meaning and convention} & \textbf{See} \\ \hline
\((X,\mu)\) & Sigma-finite measure space in the kernel model. & \hyperref[finite-kernel-model]{0.2} \\
\(H\), \(L^2(X,\mu;\mathbb C)\) & Complex square-integrable functions, identified when equal almost everywhere. & \hyperref[finite-kernel-model]{0.2} \\
$\langle f,g\rangle_{L^2(X,\mu)}$ & Inner product $\int_X f\overline g\,d\mu$ on the complex space $H$, linear in the first argument. & \hyperref[finite-kernel-model]{0.2} \\
$\|f\|_{L^2(X,\mu)}$ & Norm $\langle f,f\rangle_{L^2(X,\mu)}^{1/2}$ on $H$; both space and measure appear in the subscript. & \hyperref[finite-kernel-model]{0.2} \\
\(u_1,\ldots,u_N\) & Fixed orthonormal family with chosen measurable representatives; \(N=0\) is allowed. & \hyperref[finite-kernel-model]{0.2} \\
\(V\), \(V^\perp\) & Span of the family and its orthogonal complement in \(H\). & \hyperref[finite-kernel-model]{0.2} \\
\(K_A(x,y)\) & Finite kernel \(\sum_{i,j}a_{ij}u_i(x)\overline{u_j(y)}\); the output variable comes first. & \hyperref[finite-kernel-operator]{0.2} \\
\(T_A\) & Operator with kernel \(K_A\); it acts as \(A\) on \(V\) and vanishes on \(V^\perp\). & \hyperref[finite-kernel-operator]{0.2} \\
\(\operatorname{tr}(T_A)\) & Trace of the finite-dimensional restriction \(T_A|_V\). & \hyperref[finite-kernel-trace-definition]{0.2} \\
\(K_A(x,x)\) & Diagonal of the specified finite kernel. An arbitrary product-almost-everywhere representative need not give the same diagonal integral. & \hyperref[kernel-diagonal-representatives]{0.2} \\
\end{longtable}

\subsection*{Finite Fourier analysis}\label{v1a-notation-fourier}

\begin{longtable}{p{0.22\linewidth}p{0.58\linewidth}p{0.08\linewidth}}
\textbf{Symbol} & \textbf{Meaning and convention} & \textbf{See} \\ \hline
\(G=\mathbb Z/N\mathbb Z\) & Finite cyclic group, with \(N\geq1\) and normalized counting measure. & \hyperref[finite-circle-frequencies]{0.3} \\
$\mu_G$ & Normalized counting measure on $G$: $\mu_G(\{x\})=1/N$ for each $x\in G$. & \hyperref[finite-circle-frequencies]{0.3} \\
\(e(t)\) & Complex exponential \(\exp(2\pi i t)\). & \hyperref[finite-circle-frequencies]{0.3} \\
\(\chi_j(x)\) & Character \(e(jx/N)\), for \(0\leq j<N\). & \hyperref[finite-circle-frequencies]{0.3} \\
\(\delta_{jk}\) & Kronecker delta: one for \(j=k\), zero otherwise. & \hyperref[finite-circle-frequencies]{0.3} \\
$\langle f,g\rangle_{L^2(G,\mu_G)}$ & Inner product $\frac1N\sum_x f(x)\overline{g(x)}$ on $L^2(G,\mu_G;\mathbb C)$, linear in the first argument. & \hyperref[finite-circle-frequencies]{0.3} \\
\(\widehat f(j)\) & Fourier coefficient \(\frac1N\sum_y f(y)e(-jy/N)\). The inversion sum has no prefactor. & \hyperref[finite-fourier-coefficients]{0.3} \\
\(C_a\) & Convolution operator: \((C_a f)(x)=\frac1N\sum_y a(x-y)f(y)\), with sums over \(G\). & \hyperref[finite-convolution-operator]{0.3} \\
\end{longtable}

\subsection*{Ordinary traces and relative pairings}\label{v1a-notation-pairings}

\begin{longtable}{p{0.22\linewidth}p{0.58\linewidth}p{0.08\linewidth}}
\textbf{Symbol} & \textbf{Meaning and convention} & \textbf{See} \\ \hline
\(\mathcal T(K_A)\) & Ordinary diagonal integral \(\int_X K_A(x,x)\,d\mu(x)=\operatorname{tr}(A)\). & \hyperref[ordinary-finite-trace]{0.4} \\
$P_p(u)$ & Scalar period $\langle u,p\rangle_{L^2(X,\mu)}$, for $p\in H$. & \hyperref[relative-finite-pairing]{0.4} \\
$\mathcal R_{p,q}(K_A)$ & Relative pairing $\langle T_Aq,p\rangle_{L^2(X,\mu)}$, for $p,q\in H$. & \hyperref[relative-finite-pairing]{0.4} \\
\(1_C\) & Indicator of a measurable set \(C\). If \(\mu(C)<\infty\), then \(P_{1_C}(u)=\int_Cu\,d\mu\). & \hyperref[finite-subset-periods]{0.4} \\
\end{longtable}

\subsection*{Weights and multiplication operators}\label{v1a-notation-weights}

\begin{longtable}{p{0.22\linewidth}p{0.58\linewidth}p{0.08\linewidth}}
\textbf{Symbol} & \textbf{Meaning and convention} & \textbf{See} \\ \hline
\(L^\infty(X,\mu;\mathbb C)\) & Essentially bounded complex functions, identified when equal almost everywhere. & \hyperref[weighted-finite-trace]{0.4} \\
\(w\), \(M_w\) & A weight \(w\in L^\infty\) and its multiplication operator \(M_wf=wf\). & \hyperref[weighted-finite-trace]{0.4} \\
\(P\) & Orthogonal projection onto \(V\); distinct from the scalar period \(P_p\). & \hyperref[weighted-finite-trace]{0.4} \\
\(S_w\) & Endomorphism \(P M_wT_A|_V\), obtained by projecting the weighted output back to \(V\). & \hyperref[finite-weighted-kernel-identity]{0.4} \\
\(\mathcal W_w(K_A)\) & Weighted diagonal integral \(\int_Xw(x)K_A(x,x)\,d\mu(x)=\operatorname{tr}(S_w)\). & \hyperref[weighted-finite-trace]{0.4} \\
\end{longtable}
